% Imporditud: tools/import_eksperimendid.py
% Allikas: Eesti füüsikaolümpiaadi eksperimendi ülesannete
%   kogu (Rael Kalda, 2023), ülesanne Ü100, lahendus L100.
\setAuthor{EFO žürii}
\setRound{lõppvoor}
\setYear{2010}
\setNumber{G E1}
\setDifficulty{3}
\setTopic{elektriahelad}
\setVahendid{puitlatile kinnitatud takistustraat, kaks ühesugust lampi, tuntud takistusega (R = 12,1 $\Omega$) takisti, patarei, mõõdulint, nihik}
\setPunktid{12}
\setAllikas{Eksperimendi kogu Ü100, lahendus lk 78}
\setLahendusseis{toores}

\prob{Traat}
Määrake takistustraadi eritakistus. Hinnake mõõtemääramatust.
\textit{Märkus:} Patereid mitte lühistada, traati mitte lõigata! Pirn kannatab patarei pinget.

\hint

\solu
Ühe lambi ühendame järjestikku tuntud takistiga ning teise lambi --- takistustraadiga. Takististustraadi puhul ühendame statsionaarselt vaid ühe kontakti traadi otspunktis; teise kontakti ühendame käsitsi, hoides juhet käega vastu takistustraati mingis punktis, mille asukohta saab piki traati libistades muuta. Mõlemad ahelad ühendame patarei klemmidele. Saavutame olukorra, kus mõlemad lambid põlevad ühe heledusega ning mõõdame selles asendis kontaktide vahele jääva takistustraadi osa pikkuse L. Sellises asendis on vastava takistustraadi osa takistus võrdne teise takisti takistusega, st R = $\rho$4L/$\pi$$d_2$, millest $\rho$ = $\pi$Rd2/4L. Arvuliseks vastuseks saame $\rho$ = 1,35 $\cdot$ 10$-$6 $\Omega \cdot$ m. Suhtelise mõõtevea leiame seosest $\Delta\rho$ = 2$\Delta$d/d + $\Delta$L/L, kus nihiku viga $\Delta$d $\approx$ 0,05 mm ning pikkuse L vea leiame mitme mõõtmistulemuse (silmaga heleduse hindamisest johtuva) hajuvuse põhjal, $\Delta$L $\approx$ 3 cm.
\probend
